\documentclass[10pt,a4paper]{amsart}
\usepackage{amsmath,amssymb,amsthm,graphicx}
\usepackage{verbatim}
\newenvironment{reference}{\comment}{\endcomment}
%\newenvironment{reference}{}{}
\newenvironment{slogan}{\comment}{\endcomment}
\newenvironment{history}{\comment}{\endcomment}

% For commutative diagrams we use Xy-pic
\usepackage[all]{xy}

% We use 2cell for 2-commutative diagrams.
\xyoption{2cell}
% Theorem environments.
%
\theoremstyle{plain}
\newtheorem{theorem}[subsection]{Theorem}
\newtheorem{proposition}[subsection]{Proposition}
\newtheorem{lemma}[subsection]{Lemma}

\theoremstyle{definition}
\newtheorem{definition}[subsection]{Definition}
\newtheorem{example}[subsection]{Example}
\newtheorem{exercise}[subsection]{Exercise}
\newtheorem{situation}[subsection]{Situation}

\theoremstyle{remark}
\newtheorem{remark}[subsection]{Remark}
\newtheorem{remarks}[subsection]{Remarks}

\numberwithin{equation}{subsection}

% Macros
%
\def\lim{\mathop{\mathrm{lim}}\nolimits}
\def\colim{\mathop{\mathrm{colim}}\nolimits}
\def\Spec{\mathop{\mathrm{Spec}}}
\def\Hom{\mathop{\mathrm{Hom}}\nolimits}
\def\Ext{\mathop{\mathrm{Ext}}\nolimits}
\def\SheafHom{\mathop{\mathcal{H}\!\mathit{om}}\nolimits}
\def\SheafExt{\mathop{\mathcal{E}\!\mathit{xt}}\nolimits}
\def\Sch{\mathit{Sch}}
\def\Mor{\mathop{\mathrm{Mor}}\nolimits}
\def\Ob{\mathop{\mathrm{Ob}}\nolimits}
\def\Sh{\mathop{\mathit{Sh}}\nolimits}
\def\NL{\mathop{N\!L}\nolimits}
\def\CH{\mathop{\mathrm{CH}}\nolimits}
\def\proetale{{pro\text{-}\acute{e}tale}}
\def\etale{{\acute{e}tale}}
\def\QCoh{\mathit{QCoh}}
\def\Ker{\mathop{\mathrm{Ker}}}
\def\Im{\mathop{\mathrm{Im}}}
\def\Coker{\mathop{\mathrm{Coker}}}
\def\Coim{\mathop{\mathrm{Coim}}}

% Boxtimes
%
\DeclareMathSymbol{\boxtimes}{\mathbin}{AMSa}{"02}

%
% Macros for moduli stacks/spaces
%
\def\QCohstack{\mathcal{QC}\!\mathit{oh}}
\def\Cohstack{\mathcal{C}\!\mathit{oh}}
\def\Spacesstack{\mathcal{S}\!\mathit{paces}}
\def\Quotfunctor{\mathrm{Quot}}
\def\Hilbfunctor{\mathrm{Hilb}}
\def\Curvesstack{\mathcal{C}\!\mathit{urves}}
\def\Polarizedstack{\mathcal{P}\!\mathit{olarized}}
\def\Complexesstack{\mathcal{C}\!\mathit{omplexes}}
% \Pic is the operator that assigns to X its picard group, usage \Pic(X)
% \Picardstack_{X/B} denotes the Picard stack of X over B
% \Picardfunctor_{X/B} denotes the Picard functor of X over B
\def\Pic{\mathop{\mathrm{Pic}}\nolimits}
\def\Picardstack{\mathcal{P}\!\mathit{ic}}
\def\Picardfunctor{\mathrm{Pic}}
\def\Deformationcategory{\mathcal{D}\!\mathit{ef}}
\usepackage{cleveref}
\crefname{originallemma}{original lemma}{original lemmas}
\Crefname{originallemma}{Original lemma}{Original lemmas}
\makeatletter
\newlabel{lemma-descending-properties-morphisms}{{041J}{source}{}{sourcepointer-1}{}}
\newlabel{lemma-descending-properties-morphisms@cref}{{[originallemma][0][]041J}{[1][0][]source}{}{}{}}
\makeatother
\crefname{originalreference}{original reference}{original references}
\Crefname{originalreference}{Original reference}{Original references}
\makeatletter
\newlabel{spaces-morphisms-lemma-finite-presentation-local}{{0410}{source}{}{sourcepointer-2}{}}
\newlabel{spaces-morphisms-lemma-finite-presentation-local@cref}{{[originalreference][0][]0410}{[1][0][]source}{}{}{}}
\makeatother
\crefname{originalreference}{original reference}{original references}
\Crefname{originalreference}{Original reference}{Original references}
\makeatletter
\newlabel{descent-lemma-descending-property-locally-finite-presentation}{{02KY}{source}{}{sourcepointer-3}{}}
\newlabel{descent-lemma-descending-property-locally-finite-presentation@cref}{{[originalreference][0][]02KY}{[1][0][]source}{}{}{}}
\makeatother
\title{Being locally of finite presentation is fpqc local on the base}
\author{The Stacks Project Authors}
\date{stacks-project:spaces-descent.tex:local-snapshot-20261012; GNU FDL 1.2 or later}
\begin{document}
\maketitle
\noindent\textit{Editorial note.} All delivered mathematical text is the author's own, copied verbatim from the pinned Stacks Project chapter file spaces-descent.tex; no statement or proof step was rewritten, reordered or removed.\par\medskip
\noindent\textit{Editorial note.} The counted span cites one Morphisms of Spaces lemma and one Descent lemma plus same-chapter items, recorded in assembly.external\_references with canonical tags as navigation pointers.\par\medskip
\noindent\textit{Allowed standard prerequisites (editorial).}\begin{itemize}
\item Framework lemma for descending properties of morphisms. Stacks Project, Descent on Spaces, Lemma lemma-descending-properties-morphisms. Supplies the criterion structure that the proof verifies.
\item Local finite presentation is \textbackslash{}'etale-local on the source. Stacks Project, Morphisms of Spaces, Lemma spaces-morphisms-lemma-finite-presentation-local. Supplies assumptions (1), (2) and the chart reductions.
\item fpqc descent of local finite presentation for schemes. Stacks Project, Descent, Lemma descent-lemma-descending-property-locally-finite-presentation. Supplies the final descent step for U \textbackslash{}to Z.
\end{itemize}\medskip
\ifdefined\hypertarget\hypertarget{sourcepointer-1}{}\fi
\noindent\textit{Source reference.} Original lemma 041J: sources/stacks-chow/tags.txt, source lines 5294--5294. Cross-chapter standard prerequisite cited by the author inside the counted span; it is used as a pointer and its own proof is not part of this item.\par\medskip
\ifdefined\hypertarget\hypertarget{sourcepointer-2}{}\fi
\noindent\textit{Source reference.} Original reference 0410: sources/stacks-chow/tags.txt, source lines 5275--5275. Cross-chapter standard prerequisite cited by the author inside the counted span; it is used as a pointer and its own proof is not part of this item.\par\medskip
\ifdefined\hypertarget\hypertarget{sourcepointer-3}{}\fi
\noindent\textit{Source reference.} Original reference 02KY: sources/stacks-chow/tags.txt, source lines 3389--3389. Cross-chapter standard prerequisite cited by the author inside the counted span; it is used as a pointer and its own proof is not part of this item.\par\medskip
\begin{lemma}
\label{lemma-descending-property-locally-finite-presentation}
The property $\mathcal{P}(f) =$``$f$ is locally of finite presentation''
is fpqc local on the base.
\end{lemma}

\begin{proof}
We will use
Lemma \ref{lemma-descending-properties-morphisms}
to prove this. Assumptions (1) and (2) of that lemma follow from
Morphisms of Spaces,
Lemma \ref{spaces-morphisms-lemma-finite-presentation-local}.
Let $Z' \to Z$ be a surjective flat morphism of affine schemes over $S$.
Let $f : X \to Z$ be a morphism of algebraic spaces, and assume
that the base change $f' : Z' \times_Z X \to Z'$ is locally of
finite presentation.
We have to show that $f$ is locally of finite presentation. Let $U$ be a scheme
and let $U \to X$ be surjective and \'etale. By
Morphisms of Spaces,
Lemma \ref{spaces-morphisms-lemma-finite-presentation-local}
again, it is enough to show that $U \to Z$ is locally of finite presentation.
Since $f'$ is locally of finite presentation, and since $Z' \times_Z U$ is a
scheme \'etale over $Z' \times_Z X$ we conclude (by the same lemma again) that
$Z' \times_Z U \to Z'$ is locally of finite presentation.
As $\{Z' \to Z\}$ is an fpqc covering we conclude that
$U \to Z$ is locally of finite presentation by
Descent,
Lemma \ref{descent-lemma-descending-property-locally-finite-presentation}
as desired.
\end{proof}

\end{document}
